\documentclass[11pt,a4paper]{article}
\usepackage[margin=2.2cm]{geometry}
\usepackage{booktabs,enumitem}
\usepackage[T1]{fontenc}
\setlength{\parskip}{0.5em}\setlength{\parindent}{0pt}
\title{Why the paper's cleaner wins with depth-only ICP\\and fails with RGB-D odometry}
\date{}
\begin{document}
\maketitle

\section*{The question}
Two cleaners (our 9$\times$9 \textbf{median}, and the \textbf{paper} by Fu et al.) are compared with two trackers (depth-only ICP, RGB-D odometry). On ICL (ATE in cm, lower is better):

\begin{center}
\begin{tabular}{lcc}\toprule
 & Depth-only ICP & RGB-D odometry\\\midrule
No cleaning & 17.07 & 23.59\\
Median (ours) & 14.94 & \textbf{0.98}\\
Paper & \textbf{0.95} & 26.07\\\bottomrule
\end{tabular}
\end{center}
Two things need explaining: (A) why the paper is so good with ICP although the median is better at removing depth error, and (B) why the paper fails with RGB-D odometry.

\section*{Part A: why the paper is so good with depth-only ICP}

\subsection*{A1. What ICP measures}
Frame-to-frame ICP estimates the \emph{motion between two frames}. It aligns the point cloud of frame $t+1$ onto frame $t$. The result depends on the \textbf{difference} between the two frames, not on how close each frame is to the true depth. A constant error that appears identically in both frames does not change the estimated motion. An error that is different in each frame does.

\subsection*{A2. The surprising numbers}
\begin{center}
\begin{tabular}{lccc}\toprule
 & Depth MAE (mm) & Within 1 cm & ICP RPE (mm per step)\\\midrule
Median & \textbf{18.8} & \textbf{44.7\%} & 11.78\\
Paper & 37.7 & 33.7\% & \textbf{2.02}\\\bottomrule
\end{tabular}
\end{center}
The paper has \emph{worse} per-pixel depth error but about 6$\times$ better ICP accuracy per step. So \textbf{depth accuracy against the truth does not predict ICP accuracy.} What matters is how consistent the errors are between consecutive frames.

\subsection*{A3. The median-ICP result is mostly noise}
On ICL the camera moves only \textbf{6.5 mm per tracked step}. The median's ICP error per step is \textbf{11.78 mm}, larger than the motion itself. The tracker is then not really tracking: its error exceeds the signal. The paper's 2.02 mm per step is a real measurement of the motion.

\subsection*{A4. What each cleaner does to the noise}
\begin{itemize}[nosep]
\item \textbf{Median (ours).} Per frame, per pixel. It takes an order statistic, so it removes spikes but does not average the ordinary noise of the good pixels (median error only goes from 13.3 mm to 12.3 mm). The sensor noise is different in every frame, so the noise left in frame $t$ and frame $t+2$ is \emph{uncorrelated}.
\item \textbf{Paper (ST-DNBL).} Averages over a window of frames $t\pm2$. Our ICL run tracks every 2nd frame, so consecutive tracked frames share 3 of their 5 input frames. The noise left in consecutive outputs is therefore \emph{largely the same}, and it cancels when ICP aligns them.
\end{itemize}
Intuition: measuring how far a picture of a wall has shifted between two photos. If both photos carry the same fixed noise pattern, the shift is measured correctly. If each photo has its own random pattern, the pattern corrupts the measurement.

\subsection*{A5. Why this hurts ICP specifically}
Our ICP (\texttt{icl\_compare.py}) is point-to-plane, coarse to fine: voxels of 4/2/1 cm, normals from PCA over 3 voxels, maximum matching distance 10/5/2 cm. Point-to-plane minimises the residual \emph{along the target normal} only. Noise that differs between frames changes the estimated normals and the point positions along them. Flat walls and floors constrain motion along the plane poorly, so a small bias at each step is not corrected. The bias repeats over 150 steps and shows up as drift.

\subsection*{A6. Why spikes do not hurt the paper under ICP}
Voxel averaging, the matching-distance cap (2 cm at the finest level) and the point-to-plane residual all work against isolated spikes. A spike is usually rejected rather than fitted. So the paper can keep its spikes and still track well.

\section*{Part B: why the paper fails with RGB-D odometry}
\begin{enumerate}[nosep]
\item The paper's bilateral filter averages only neighbours with \emph{similar} depth. A spike has no similar neighbours, so it averages with itself and survives. The paper's RMSE stays at 186 mm, about the same as the uncleaned input (187 mm).
\item Hybrid RGB-D odometry works at full resolution. It uses the depth at every pixel, with no voxel averaging. It does ignore pixel pairs whose depth differs by more than 3 cm (Open3D's \texttt{depth\_diff\_max}, the default our code uses), but a spike smaller than that still enters the alignment.
\item The surviving spikes pull the alignment off. With the paper's output, 49 of 149 steps were off by more than 20 mm; with the median, none. \emph{The size of the paper's remaining spikes, and so exactly how they get past the 3 cm gate, has not been checked.}
\end{enumerate}
The median removes the spikes, which is the one thing RGB-D odometry needs, so it reaches 0.98 cm (perfect depth gives 0.81 cm).

\section*{Summary}
\begin{center}
\begin{tabular}{p{3.2cm}p{5.6cm}p{5.6cm}}\toprule
 & Median & Paper\\\midrule
Removes spikes & yes & no\\
Averages ordinary noise & no & yes, over $t\pm2$ frames\\
Noise similar between consecutive frames & no & yes (hypothesis)\\
Depth-only ICP needs & consistency between frames $\to$ \textbf{fails} & consistency between frames $\to$ \textbf{works}\\
RGB-D odometry needs & no spikes $\to$ \textbf{works} & no spikes $\to$ \textbf{fails}\\\bottomrule
\end{tabular}
\end{center}

\section*{What is measured and what is a hypothesis}
\textbf{Measured:} all numbers in the tables above, and the fact that the paper has worse depth error but better ICP accuracy.

\textbf{Hypothesis, not yet tested:} that the paper wins because its temporal averaging makes consecutive frames share their noise. (The README gives a related explanation, ``smooth surfaces give good normals''.) Three experiments would test it:
\begin{enumerate}[nosep]
\item Run the paper's pipeline with the temporal window set to zero (spatial bilateral only). If ICP drift rises towards the median's 15 cm, temporal averaging is the cause.
\item Apply a 5-frame temporal median, then the 9$\times$9 median. If this reaches about 1 cm with ICP, the same conclusion holds.
\item Measure the error of the \emph{difference between consecutive cleaned frames} against the true difference. The hypothesis predicts a much lower value for the paper, even though its per-frame MAE is higher.
\end{enumerate}
\end{document}
